% Source context: Garre-Rubio--Lootens--Molnar, arXiv:2203.12563v3,
% REsubmission.tex, algcond, defsym, Section 5, and Appendix B.
% The physical adjoint hypothesis below is additional to algebraic closedness.
% Source-only declarations: compiler and axiom-guard validation are pending.
% No statement or proof carries a checked marker.

\section{Normal blocks and arbitrary-boundary adjoints}
\label{sec:glm_normal_adjoint_boundary}

For an operator tensor $T$ of bond dimension $\chi$, write
$O^N_{T,X}(\sigma,\tau)=\tr(XT^{\sigma_1\tau_1}\cdots
    T^{\sigma_N\tau_N})$ and $O_T^N=O^N_{T,I}$.
The physical adjoint tensor is
$(T^\sharp)^{ij}=\overline{T^{ji}}$, where the bar denotes entrywise
complex conjugation. Thus
$(T^\sharp)^{ij}_{\alpha\beta}=\overline{T^{ji}_{\alpha\beta}}$:
the physical indices are exchanged, while the virtual indices retain
their order. At $N=0$ the contraction is the scalar $\tr X$.

The algebraic arbitrary-boundary closedness of
\cite[Section~2]{GarreRubio2023MPOSym} does not supply a physical
adjoint operation. The periodic adjoint identity below is an explicit
additional hypothesis, consistent with the extra physical symmetry
structure considered in \cite[Section~5]{GarreRubio2023MPOSym}.

\subsection{From periodic duality to exact similarities}

\begin{theorem}[Physical adjunction with an arbitrary boundary]
    \label{thm:glm_boundary_physical_adjoint}
    \lean{MPOTensor.conjTranspose_mpoWithBoundary}
    For every $X\in\MN{\chi}$ and $N\geq0$,
    \begin{align}
        (O^N_{T,X})^\dagger & =O^N_{T^\sharp,\overline X}.
        \label{eq:glm_boundary_physical_adjoint}
    \end{align}
\end{theorem}
\begin{proof}
    The $(\sigma,\tau)$ entry of the left side is
    $\overline{\tr(XT^{\tau_1\sigma_1}\cdots
            T^{\tau_N\sigma_N})}$.
    Entrywise conjugation preserves matrix multiplication, its order,
    and the trace, giving the right side. In particular the boundary
    is $\overline X$, without a virtual transpose.
\end{proof}

\begin{theorem}[Boundary transport through a two-sided similarity]
    \label{thm:glm_boundary_similarity}
    \lean{MPOTensor.mpoWithBoundary_eq_of_conj}
    Let $P,Q$ have bond dimensions $\chi,\eta$, respectively, and let
    $V:\C^\chi\to\C^\eta$, $W:\C^\eta\to\C^\chi$ satisfy
    $VW=I_\eta$, $WV=I_\chi$, and $VP^{ij}W=Q^{ij}$ for every $i,j$.
    Then $O^N_{P,X}=O^N_{Q,VXW}$ for every $X$ and $N\geq0$.
\end{theorem}
\begin{proof}
    Cancelling successive factors $WV$ extends the letter identity to
    every word. Equivalently, $P^w=WQ^wV$, including the empty word
    because $WV=I_\chi$. Cyclicity of the rectangular trace gives
    $\tr(XWQ^wV)=\tr(VXWQ^w)$.
\end{proof}

\begin{definition}[Periodic physical adjoint duality]
    \label{def:glm_periodic_adjoint_family}
    \lean{MPOTensor.IsPeriodicAdjointFamily,
        MPOTensor.IsPeriodicAdjointFamily.sameMPV₂Pos}
    Let $T_a$ be operator tensors with bond dimensions $\chi_a$, and
    let $a\mapsto a^\vee$ be a map on their labels. Periodic physical
    adjoint duality means that, for every label $a$ and every $N>0$,
    \begin{align}
        O^N_{T_{a^\vee}} & =(O^N_{T_a})^\dagger.
        \label{eq:glm_periodic_adjoint}
    \end{align}
    Equivalently, regarding the physical pair $(i,j)$ as one index,
    $T_a^\sharp$ and $T_{a^\vee}$ have equal periodic matrix product
    vectors at every positive length.
\end{definition}

\begin{theorem}[Equality of dual bond dimensions]
    \label{thm:glm_dual_bond_dimension}
    \lean{MPOTensor.IsPeriodicAdjointFamily.bondDim_dual}
    \uses{def:glm_periodic_adjoint_family,def:normal}
    Suppose (\ref{eq:glm_periodic_adjoint}) holds,
    $(a^\vee)^\vee=a$, and every $T_a$, on its paired physical alphabet,
    is normal with $\chi_a>0$. Then $\chi_a=\chi_{a^\vee}$ for every $a$.
    Here normality means that some positive blocking length makes the
    word matrices span $\MN{\chi_a}$.
\end{theorem}
\begin{proof}
    \uses{thm:mps_equality_reduction_exists}
    Apply the normal-target reduction theorem to the equal periodic
    families $T_a^\sharp$ and $T_{a^\vee}$.
    It gives $V_aW_a=I_{\chi_{a^\vee}}$, hence
    $\chi_{a^\vee}\leq\chi_a$. Applying the same theorem at $a^\vee$
    gives $\chi_{(a^\vee)^\vee}\leq\chi_{a^\vee}$.
    Involutivity proves the reverse inequality.
    This uses the equality case of
    \cite[Proposition~20]{Molnar2018SemiInjective}; normality is required
    for the target in each application.
\end{proof}

\begin{theorem}[Exact dual gauges and arbitrary-boundary adjoints]
    \label{thm:glm_normal_adjoint_gauge}
    \lean{MPOTensor.IsPeriodicAdjointFamily.exists_boundaryGauge}
    \uses{def:glm_periodic_adjoint_family,def:normal}
    Under the hypotheses of Theorem~\ref{thm:glm_dual_bond_dimension},
    for each $a$ there exist maps
    $V_a:\C^{\chi_a}\to\C^{\chi_{a^\vee}}$ and
    $W_a:\C^{\chi_{a^\vee}}\to\C^{\chi_a}$ with
    \begin{align}
        V_aW_a            & =I_{\chi_{a^\vee}},     & W_aV_a & =I_{\chi_a},
        \label{eq:glm_dual_inverses}                                        \\
        (T_a^\sharp)^{ij} & =W_aT_{a^\vee}^{ij}V_a.
        \label{eq:glm_dual_letters}
    \end{align}
    The same maps satisfy, for every $X\in\MN{\chi_a}$ and $N\geq0$,
    \begin{align}
        (O^N_{T_a,X})^\dagger
         & =O^N_{T_{a^\vee},\,V_a\overline XW_a}.
        \label{eq:glm_dual_boundary}
    \end{align}
    No one-site injectivity, canonical normalization, or inequivalence
    between distinct blocks is required.
\end{theorem}
\begin{proof}
    \uses{thm:mps_equality_reduction_exists,thm:glm_dual_bond_dimension,
        thm:glm_boundary_physical_adjoint,thm:glm_boundary_similarity}
    A normal-target reduction gives $V_aW_a=I$ and
    $V_a(T_a^\sharp)^{ij}W_a=T_{a^\vee}^{ij}$.
    Equality of the two dimensions makes the right inverse two-sided,
    so $W_aV_a=I$. Multiplying the letter identity by $W_a$ and $V_a$
    proves (\ref{eq:glm_dual_letters}). Combine
    (\ref{eq:glm_boundary_physical_adjoint}) with
    Theorem~\ref{thm:glm_boundary_similarity} to obtain
    (\ref{eq:glm_dual_boundary}). The two-sided identities also cover
    $N=0$, although the periodic hypothesis only concerns $N>0$.
\end{proof}

In the following diagrams, a west virtual leg is a matrix row and an
east virtual leg is a matrix column. The north physical leg $i$ is
an output, and the south physical leg $j$ is an input. The first
identity contracts $p,q$ of dimension $\chi_{a^\vee}$; the second
contracts $\alpha,\beta$ of dimension $\chi_a$ and writes
$B=V_a\overline XW_a$.
% TENKZ-NORMAL-ADJOINT-LETTERS-BEGIN
% Formula: (T_a^sharp)^{ij}_{alpha beta}
% = sum_{p,q} (W_a)_{alpha p} (T_{a^vee}^{ij})_{pq} (V_a)_{q beta}.
% Physical i (north output) and j (south input) range over Fin d.
% External alpha,beta range over Fin chi_a; contracted p,q range over
% Fin chi_{a^vee}. West ports are matrix rows; east ports are columns.
% There is no conjugation or transpose on W_a and V_a.
% Boundary signature: open:e, open:w, phys:n, phys:s.
\begin{tenkzequation}
    \begin{tenkz}[rows={wire}, cols=1, bonds=none]
        \tn[name=S, skin=box,size=l,
            ports={180:virtual:$\alpha$,0:virtual:$\beta$,
                    90:physical:$i$,270:physical:$j$}]{T_a^\sharp}
    \end{tenkz}
    $=$
    \begin{tenkz}[rows={wire}, cols=3, bonds=none]
        \tn[at=(1,1),name=W,skin=box,size=l,
            ports={180:virtual:$\alpha$,0:virtual}]{W_a}
        \tn[at=(1,2),name=T,skin=box,size=l,
            ports={180:virtual,0:virtual,
                    90:physical:$i$,270:physical:$j$}]{T_{a^\vee}}
        \tn[at=(1,3),name=V,skin=box,size=l,
            ports={180:virtual,0:virtual:$\beta$}]{V_a}
        \tnwire{W.0}{T.180}
        \tnwire{T.0}{V.180}
    \end{tenkz}
\end{tenkzequation}
% TENKZ-NORMAL-ADJOINT-LETTERS-END
% TENKZ-NORMAL-ADJOINT-BOUNDARY-BEGIN
% Formula: B_{pq} = sum_{alpha,beta}
% (V_a)_{p alpha} conjugate(X_{alpha beta}) (W_a)_{beta q}.
% External p,q range over Fin chi_{a^vee}; contracted alpha,beta range
% over Fin chi_a. Conjugation is entrywise, with no virtual transpose.
% West ports are matrix rows; east ports are columns. No physical legs.
% Boundary signature: open:e, open:w.
\begin{tenkzequation}
    \begin{tenkz}[rows={wire}, cols=1, bonds=none]
        \tn[name=B,skin=box,size=l,
            ports={180:virtual:$p$,0:virtual:$q$}]{B}
    \end{tenkz}
    $=$
    \begin{tenkz}[rows={wire}, cols=3, bonds=none]
        \tn[at=(1,1),name=V,skin=box,size=l,
            ports={180:virtual:$p$,0:virtual}]{V_a}
        \tn[at=(1,2),name=X,skin=box,size=l,
            ports={180:virtual,0:virtual}]{\overline X}
        \tn[at=(1,3),name=W,skin=box,size=l,
            ports={180:virtual,0:virtual:$q$}]{W_a}
        \tnwire{V.0}{X.180}
        \tnwire{X.0}{W.180}
    \end{tenkz}
\end{tenkzequation}
% TENKZ-NORMAL-ADJOINT-BOUNDARY-END

\subsection{Literal block sums and canonical parents}

Let the labels be $a=0,\ldots,r-1$, put $\chi=\sum_a\chi_a$, and take
the literal unweighted block sum $T^{ij}=\bigoplus_aT_a^{ij}$.
Write $J_a:\C^{\chi_a}\to\bigoplus_b\C^{\chi_b}$ for its coordinate
inclusions and $X_{aa}=J_a^\dagger XJ_a$ for the diagonal blocks of an
arbitrary $X\in\MN{\chi}$.

\begin{theorem}[Diagonal boundaries of a block sum]
    \label{thm:glm_adjoint_block_diagonal}
    \lean{MPOTensor.mpoWithBoundary_blockSum}
    At every length $N\geq0$,
    \begin{align}
        O^N_{T,X} & =\sum_aO^N_{T_a,X_{aa}}.
        \label{eq:glm_adjoint_block_diagonal}
    \end{align}
\end{theorem}
\begin{proof}
    Each word in $T$ is block diagonal with blocks the corresponding
    words in $T_a$. Only the diagonal blocks of $X$ contribute to its
    trace pairing. The empty word is the direct-sum identity, so the
    same argument applies at $N=0$.
\end{proof}

\begin{definition}[An arbitrary boundary in one summand]
    \label{def:glm_adjoint_block_inclusion}
    \lean{MPOTensor.embedBlockBoundary,
        MPOTensor.mpoWithBoundary_blockSum_embedBlockBoundary}
    For $Z\in\MN{\chi_a}$ define $\iota_a(Z)=J_aZJ_a^\dagger$.
    Then $O^N_{T,\iota_a(Z)}=O^N_{T_a,Z}$ for every $N\geq0$.
    This inserts a full matrix boundary, rather than only a scalar
    multiple of the identity in that summand.
\end{definition}

\begin{definition}[Adjoint transport on the whole boundary space]
    \label{def:glm_block_adjoint_transport}
    \lean{MPOTensor.blockSumAdjointBoundary}
    \uses{def:glm_adjoint_block_inclusion}
    Given dual labels and maps $V_a,W_a$ of the dimensions in
    Theorem~\ref{thm:glm_normal_adjoint_gauge}, define
    \begin{align}
        \mathcal B(X)
         & =\sum_a\iota_{a^\vee}
        \!\left(V_a\overline{X_{aa}}W_a\right).
        \label{eq:glm_block_adjoint_transport}
    \end{align}
    This conjugate-linear boundary map is independent of the chain length.
\end{definition}

\begin{definition}[Length-independent boundary adjoint closure]
    \label{def:glm_boundary_adjoint_closed}
    \lean{MPOTensor.IsBoundaryAdjointClosed}
    A tensor $T$ is boundary adjoint closed if, for each boundary $X$,
    there exists one boundary $Y$ such that
    $(O^N_{T,X})^\dagger=O^N_{T,Y}$ for every $N>0$.
    The output boundary is chosen before the length.
\end{definition}

\begin{theorem}[Adjoint closure from normal periodic dual blocks]
    \label{thm:glm_normal_block_adjoint}
    \lean{MPOTensor.IsPeriodicAdjointFamily.exists_blockSumAdjointBoundary,
        MPOTensor.IsPeriodicAdjointFamily.isBoundaryAdjointClosed_blockSum}
    \uses{def:glm_periodic_adjoint_family,def:normal,
        def:glm_block_adjoint_transport,def:glm_boundary_adjoint_closed}
    Suppose $a\mapsto a^\vee$ is involutive, the $T_a$ are normal with
    positive bond dimensions, and (\ref{eq:glm_periodic_adjoint}) holds.
    There exist maps $V_a,W_a$ satisfying
    (\ref{eq:glm_dual_inverses}) and (\ref{eq:glm_dual_letters}) for every
    $a$, such that every $X$ and $N\geq0$ satisfy
    \begin{align}
        (O^N_{T,X})^\dagger & =O^N_{T,\mathcal B(X)}.
        \label{eq:glm_normal_block_adjoint}
    \end{align}
    In particular the literal block sum is boundary adjoint closed.
\end{theorem}
\begin{proof}
    \uses{thm:glm_normal_adjoint_gauge,thm:glm_adjoint_block_diagonal,
        def:glm_adjoint_block_inclusion,def:glm_block_adjoint_transport}
    Choose the exact dual gauges block by block. Taking adjoints in
    (\ref{eq:glm_adjoint_block_diagonal}), then applying
    (\ref{eq:glm_dual_boundary}), gives
    \begin{align*}
        (O^N_{T,X})^\dagger
         & =\sum_aO^N_{T_{a^\vee},V_a\overline{X_{aa}}W_a}
        =\sum_aO^N_{T,\iota_{a^\vee}(V_a\overline{X_{aa}}W_a)}
        =O^N_{T,\mathcal B(X)}.
    \end{align*}
    The last equality is linearity in the boundary. Off-diagonal blocks
    of the input boundary contribute zero; no separate adjoint closure
    of an individual label is assumed.
\end{proof}

\begin{theorem}[Canonical parent symmetry from normal adjoint duality]
    \label{thm:glm_normal_adjoint_parent}
    \lean{MPOTensor.IsBoundaryCompatible.parentHamiltonianES_commute_mpo_block_of_periodicAdjoint}
    \uses{def:glm_boundary_closedness,def:glm_periodic_adjoint_family,def:normal}
    Under the hypotheses of Theorem~\ref{thm:glm_normal_block_adjoint},
    let the literal block sum $T$ be boundary compatible with a state
    tensor $A$. Let $P_{A,L}$ be the orthogonal projection onto its
    length-$L$ arbitrary-boundary state space and let $H_{A,L,N}$ be the
    periodic canonical parent formed by all cyclic translates of
    $I-P_{A,L}$. For every $0<L\leq N$ and every label $a$,
    \begin{align}
        [H_{A,L,N},O^N_{T_a}] & =0.
        \label{eq:glm_normal_adjoint_parent}
    \end{align}
    No normality or injectivity of $A$ is required.
\end{theorem}
\begin{proof}
    \uses{thm:glm_normal_block_adjoint,thm:glm_boundary_parent_blocks}
    Theorem~\ref{thm:glm_normal_block_adjoint} gives adjoint closure of
    the entire boundary range at length $L$. Compatibility makes
    $\mathcal S_A^L$ invariant under that range. Its orthogonal
    complement is invariant as well, since
    $\langle u,Bv\rangle=\langle B^\dagger u,v\rangle=0$ for
    $u\in\mathcal S_A^L$ and $v\in(\mathcal S_A^L)^\perp$.
    Thus $I-P_{A,L}$ commutes with each local boundary operator.
    The summand selector $J_aJ_a^\dagger$ commutes with every letter of
    $T$ and extracts $O^N_{T_a}$. Cutting at each cyclic window, including
    wrapping windows, therefore gives commutation with every translated
    local parent term. Theorem~\ref{thm:glm_boundary_parent_blocks}
    gives this commutation criterion for the literal block sum.
    Applying it with the derived adjoint closure proves
    (\ref{eq:glm_normal_adjoint_parent}).
\end{proof}

These results give a sufficient physical $*$-family criterion for
canonical-parent commutation. They neither construct a $C^*$-weak Hopf
algebra, its unit, or its canonical integral, nor identify the canonical
interaction with a weak-Hopf average. In the notation of
\cite[Appendix~B]{GarreRubio2023MPOSym}, a normalized average with
$E(I)=Q$ satisfies $E(I-P)=Q-E(P)$, whereas the ambient complement is
$I-E(P)$. The two coincide only when $Q=I$. No nonzeroness or spectral
gap assertion for either interaction is made here.

\subsection{A derived algebraic counterexample}
\label{sec:glm_adjoint_counterexample}

The following example is a direct algebraic construction. It shows why
normality, exact boundary multiplication, and compatibility do not replace
(\ref{eq:glm_periodic_adjoint}). It is not a counterexample to the
physical $*$-representation or weak-Hopf assumptions discussed in
\cite[Section~5]{GarreRubio2023MPOSym}.

\begin{definition}[A nonorthogonal idempotent and a compatible product state]
    \label{def:glm_adjoint_counterexample}
    \lean{MPOTensor.BoundaryAdjointCounterexample.tensor,
        MPOTensor.BoundaryAdjointCounterexample.state,
        MPOTensor.BoundaryAdjointCounterexample.boundary_one_apply,
        MPOTensor.BoundaryAdjointCounterexample.state_boundary_one_apply}
    Take physical dimension $2$, bond dimension $1$, and scalar tensor
    letters $T^{ij}=P_{ij}$ and $A^i=\delta_{i0}$, where
    \begin{align}
        P   & =\begin{pmatrix}1&1\\0&0\end{pmatrix}, &
        A^0 & =1,                                    & A^1 & =0.
        \label{eq:glm_counterexample_tensors}
    \end{align}
    Identify each bond boundary with its scalar entry. At one site,
    $O^1_{T,x}=xP$ and $\psi^1_{A,y}=y\ket{0}$.
\end{definition}

\begin{theorem}[Injectivity of both bond-one tensors]
    \label{thm:glm_counterexample_injective}
    \lean{MPOTensor.BoundaryAdjointCounterexample.tensor_isInjective,
        MPOTensor.BoundaryAdjointCounterexample.state_isInjective,
        MPOTensor.BoundaryAdjointCounterexample.tensor_isNormal,
        MPOTensor.BoundaryAdjointCounterexample.state_isNormal}
    \uses{def:glm_adjoint_counterexample,def:injective,def:normal}
    Both $T$, on its paired physical alphabet, and $A$ are one-site
    injective, and hence normal.
\end{theorem}
\begin{proof}
    Each family contains the scalar matrix $1$, which spans
    $\MN{1}=\C$. Thus blocking length one suffices.
    This is normality of the tensor letters; it is not normality of
    the physical matrix $P$ as an operator on $\C^2$.
\end{proof}

\begin{theorem}[Exact fusion and action identities]
    \label{thm:glm_counterexample_local}
    \lean{MPOTensor.BoundaryAdjointCounterexample.mulTensor_tensor,
        MPOTensor.BoundaryAdjointCounterexample.actTensor_state}
    \uses{def:glm_adjoint_counterexample}
    The stacked operator tensor is exactly $T$, and its action on
    the state tensor is exactly $A$:
    \begin{align}
        \sum_{k=0}^1T^{ik}T^{kj} & =T^{ij}, &
        \sum_{j=0}^1T^{ij}A^j    & =A^i.
        \label{eq:glm_counterexample_local}
    \end{align}
\end{theorem}
\begin{proof}
    Direct multiplication gives $P^2=P$ and $P\ket{0}=\ket{0}$.
    The tensor products of the one-dimensional bond spaces are
    canonically one-dimensional, so these are the stated letter identities.
\end{proof}

\begin{theorem}[Algebraic closedness and compatibility at every length]
    \label{thm:glm_counterexample_closed}
    \lean{MPOTensor.BoundaryAdjointCounterexample.tensor_isBoundaryClosed,
        MPOTensor.BoundaryAdjointCounterexample.tensor_isBoundaryCompatible,
        MPOTensor.BoundaryAdjointCounterexample.mem_commutingBoundaryAlgebra}
    \uses{def:glm_adjoint_counterexample,def:glm_boundary_closedness}
    The tensor $T$ is boundary closed and boundary compatible with $A$.
    For all scalar boundaries $x,y$ and every $N>0$, the same output
    boundary $xy$ satisfies
    \begin{align}
        O^N_{T,x}O^N_{T,y}    & =O^N_{T,xy},    &
        O^N_{T,x}\psi^N_{A,y} & =\psi^N_{A,xy}.
        \label{eq:glm_counterexample_boundary}
    \end{align}
    Every boundary commutes with every virtual tensor letter.
\end{theorem}
\begin{proof}
    \uses{thm:glm_counterexample_local,
        thm:glm_boundary_closedness_of_decomposition}
    The exact identities (\ref{eq:glm_counterexample_local}) give
    one-copy biorthogonal fusion and action decompositions with both
    analysis and synthesis maps equal to $1$.
    Boundary transport therefore multiplies the two scalar boundaries.
    Explicitly, contraction gives $O^N_{T,x}=xP^{\otimes N}$ and
    $\psi^N_{A,y}=y\ket{0}^{\otimes N}$; idempotence and the fixed-vector
    identity give (\ref{eq:glm_counterexample_boundary}).
    All virtual letters and boundaries are scalars, so they commute.
\end{proof}

\begin{theorem}[Failure of adjoint closure of the whole boundary range]
    \label{thm:glm_counterexample_not_adjoint}
    \lean{MPOTensor.BoundaryAdjointCounterexample.adjoint_not_mem_boundaryRange,
        MPOTensor.BoundaryAdjointCounterexample.not_boundaryAdjointClosed}
    \uses{def:glm_adjoint_counterexample}
    Already at length one,
    \begin{align}
        (O^1_{T,1})^\dagger=P^\dagger
         & =\begin{pmatrix}1&0\\1&0\end{pmatrix}
        \notin\{O^1_{T,y}:y\in\C\}=\C P.
        \label{eq:glm_counterexample_not_adjoint}
    \end{align}
    Thus even the full one-site boundary range is not adjoint closed.
\end{theorem}
\begin{proof}
    Every matrix $yP$ has $(1,0)$ entry zero, while
    $(P^\dagger)_{10}=1$. No output boundary can represent this adjoint,
    even if it is allowed to depend on the length.
\end{proof}

The canonical one-site parent also exhibits the obstruction directly.
Since $\mathcal S_A^1=\C\ket{0}$, its orthogonal complementary projection is
\begin{align}
    h=I-\ketbra{0}{0} & =\begin{pmatrix}0&0\\0&1\end{pmatrix}, &
    hP                & =0,                                    & Ph & =\begin{pmatrix}0&1\\0&0\end{pmatrix}.
    \label{eq:glm_counterexample_parent}
\end{align}
Consequently $[h,P]\ne0$, despite exact arbitrary-boundary compatibility
and commutation of all virtual boundaries with the tensor letters.
The missing property is physical adjoint closure: on the singleton
label set, periodic physical adjoint duality would require
$P^\dagger=P$ at length one, which fails.
